\scanpage{19}
% archive sheet number: 18

\uline{Remarque} On n'a jamais, dans le calcul précédent, utilisé la structure additive des $A(X)$, seulement leur structure de monoïdes. Donc il y a lieu de formuler la théorie des correspondances en prenant \deleted{des} $X\mapsto A(X)$ \deleted{qui sont des} \inserted{foncteur dans les} ``monoïdes bivariants'' \ldots\,. \deleted{Il faudrait être \deleted{\unsure{précis|prudent}} \inserted{soigneux} pour les signes, dans le cas de monoïdes non commutatifs, en prenant des monoïdes gradués [les degrés en \unsure{seront additifs|seront astreints|anticommutatifs}]} Ce point de vue cependant ne semble pas très commode dans le cas où les 2 formules de projections \inserted{(gauche et droite)} ne sont pas valables l'une et l'autre, et qu'il faut introduire des signes. Il faudrait alors remplacer les monoïdes par des ``monoïdes jaugés'' $(A^i)_{i\in\unsure{\mathbb Z|Z|\mathbf{Z}}}$, loi de composition $A^i\times A^{\unsure{j|i}}\mathbin{\unsure{-|\to}} A^{i+j}$ (NB : L'ensemble somme des $A^i$ \deleted{\unsure{sera|est}} muni d'une structure de monoïde) et \deleted{\unsure{autom.} involutif d'ordre} \inserted{application bijective} $\sigma^i\colon A^i\mathbin{-} A^i$ satisfaisant $ {\sigma^i}^{2}=\mathrm{id}$, $\sigma^i=\mathrm{id}$ si $i$ pair, \deleted{pour} enfin chaque $A^*$ muni d'un \uline{degré} \inserted{$d(A^*)$} qui permet d'imposer le degré des $f_*$ des $(f_*,f^*)\colon A^{\unsure{*|\bullet}}\mathbin{-}B^*$ comme \deleted{\unsure{f}} $d(B^*)-d(A^{\unsure{*|\circ}})$. On peut alors \unsure{exprimer|imposer|imaginer} comment les signes doivent être calculés dans la formule de projection gauche (si on veut que la droite soit valable).
